\documentclass[11pt,a4paper]{article}
\usepackage[utf8]{inputenc}
\usepackage[T1]{fontenc}
\usepackage[spanish,provide=*]{babel}
\addto\captionsspanish{\renewcommand{\tablename}{Tabla}}
\usepackage{lmodern}
\usepackage{microtype}
\usepackage[margin=2.5cm]{geometry}
\usepackage{amsmath,amssymb}
\usepackage{graphicx}
\usepackage{booktabs}
\usepackage{adjustbox}
\usepackage{pdflscape}
\usepackage{float}
\usepackage[font=small,labelfont=bf]{caption}
\usepackage{enumitem}
\usepackage{xcolor}
\usepackage{listings}
\usepackage[round]{natbib}
\usepackage[hidelinks]{hyperref}
\setlength{\parskip}{0.5em}
\setlength{\parindent}{0pt}
\setlist{itemsep=0.2em}

% Texto en rojo = algo que hay que revisar/actualizar después de volver a correr el do-file.
% Para sacar el color: \renewcommand{\revisar}[1]{#1}
\newcommand{\revisar}[1]{\textcolor{red}{#1}}
% Inserta una tabla exportada desde Stata (se achica si no entra)
\newcommand{\tabla}[1]{\begin{adjustbox}{max width=\linewidth, max totalheight=0.72\textheight}\input{#1}\end{adjustbox}}
% Figura que todavía no existe (se genera al correr el do-file): muestra un recuadro
\newcommand{\figura}[2][0.75\linewidth]{\IfFileExists{#2}{\includegraphics[width=#1]{#2}}{\fbox{\parbox{\dimexpr#1-2\fboxsep-2\fboxrule\relax}{\centering\vspace{2em}\textcolor{red}{Figura pendiente: se genera al correr el do-file}\\[0.5em]\texttt{\detokenize{#2}}\vspace{2em}}}}}
% Notas al pie de tablas
\newcommand{\notas}[1]{\par\vspace{0.3em}\begin{minipage}{0.92\linewidth}\footnotesize\textit{Notas:} #1\end{minipage}}
% Consignas en cursiva
\newcommand{\consigna}[1]{\par\textit{#1}\par}

\lstset{
  language={},
  basicstyle=\ttfamily\scriptsize,
  breaklines=true,
  columns=fullflexible,
  keepspaces=true,
  showstringspaces=false,
  frame=single,
  rulecolor=\color{gray!50},
  commentstyle=\color{gray},
  morecomment=[l]{*\ },
  morecomment=[l]{//},
  morecomment=[s]{/*}{*/},
  tabsize=4,
  inputencoding=utf8,
  extendedchars=true,
  literate={á}{{\'a}}1 {é}{{\'e}}1 {í}{{\'i}}1 {ó}{{\'o}}1 {ú}{{\'u}}1
           {Á}{{\'A}}1 {É}{{\'E}}1 {Í}{{\'I}}1 {Ó}{{\'O}}1 {Ú}{{\'U}}1
           {ñ}{{\~n}}1 {Ñ}{{\~N}}1 {ü}{{\"u}}1 {Ü}{{\"U}}1
           {¿}{{?`}}1 {¡}{{!`}}1 {’}{{'}}1 {‘}{{`}}1 {“}{{``}}1 {”}{{''}}1
           {ã}{{\~a}}1 {õ}{{\~o}}1 {â}{{\^a}}1 {ê}{{\^e}}1 {ô}{{\^o}}1 {à}{{\`a}}1
           {Ã}{{\~A}}1 {Õ}{{\~O}}1 {Â}{{\^A}}1 {Ê}{{\^E}}1 {Ô}{{\^O}}1 {Ç}{{\c C}}1 {×}{{$\times$}}1
           {°}{{$^\circ$}}1 {–}{{--}}1 {—}{{---}}1 {·}{{$\cdot$}}1 {ç}{{\c c}}1
}
\begin{document}

\begin{center}
{\Large\bfseries Problem Set 2: Multiple hypothesis testing}\\[0.5em]
E178: Economía Aplicada --- Universidad de San Andrés\\
Fecha de entrega: 30 de agosto de 2024\\[0.5em]
\textbf{Integrantes:} Juan Batalla, Nina Di Costanzo Pereira, Catalina Gallardo e Ignacio Yenarópulos
\end{center}

En este problem set corregimos la inferencia de la Tabla 2 de \citet{attanasio}, que estudia con un experimento aleatorio en Colombia los efectos de un programa de visitas domiciliarias de estimulación psicosocial sobre las habilidades cognitivas y socioemocionales de los niños y sobre las inversiones de los padres (materiales y de tiempo).

\section*{1. Replicación de la Tabla 2}

Para cada outcome estimamos por MCO una regresión de la medida en el seguimiento contra la variable de tratamiento, su valor de base (o el control de base indicado en el do-file provisto), el sexo del niño y efectos fijos de evaluador, con errores estándar clusterizados a nivel de municipio (96 clusters). La Tabla~\ref{tab:rep} presenta los resultados, con un panel por cada conjunto de outcomes, los outcomes en las columnas y el coeficiente de \texttt{treat} en las filas.

\begin{table}[H]\centering
\caption{Efectos del tratamiento (replicación de la Tabla 2 de Attanasio et al., 2020)}\label{tab:rep}
\small
\textbf{Panel A: habilidades cognitivas del niño}\\[0.2em]
\tabla{"../output/Tabla 1 - Panel A.tex"}\\[0.8em]
\textbf{Panel B: habilidades socioemocionales del niño}\\[0.2em]
\tabla{"../output/Tabla 1 - Panel B.tex"}\\[0.8em]
\textbf{Panel C: inversiones materiales}\\[0.2em]
\tabla{"../output/Tabla 1 - Panel C.tex"}\\[0.8em]
\textbf{Panel D: inversiones de tiempo}\\[0.2em]
\tabla{"../output/Tabla 1 - Panel D.tex"}
\notas{Cada columna es una regresión distinta. Todas las medidas están estandarizadas; las marcadas con (-) están invertidas, de modo que un valor más alto indica mejor comportamiento. Panel A: Bayley (columnas 1--4) y MacArthur (5--6). Panel B: ICQ (1--3) y ECBQ (4--5). Paneles C y D: ítems del FCI. Controles: valor de base, sexo y efectos fijos de evaluador. Errores estándar clusterizados por municipio entre paréntesis. * $p<0.10$, ** $p<0.05$, *** $p<0.01$.}
\end{table}

Nuestros resultados son muy parecidos a los del artículo (Tabla~\ref{tab:paper}). Las diferencias en los coeficientes son pequeñas (la mayor es en \textit{Read to child}: 0.403 contra 0.362), los errores estándar son prácticamente idénticos y los tamaños de muestra difieren en una observación. Las diferencias pueden deberse, como advierte la consigna, a que no contamos exactamente con las mismas variables que los autores. Sin corregir por múltiples hipótesis, el tratamiento tiene efectos positivos y significativos sobre el desarrollo cognitivo (Bayley cognitivo y lenguaje receptivo), sobre varias inversiones materiales (tipos de materiales de juego y juguetes para aprender formas, con un efecto negativo sobre los libros para colorear) y sobre las cinco medidas de inversión de tiempo. Ninguno de los ítems socioemocionales es significativo individualmente.

\begin{table}[H]\centering
\caption{Comparación con los resultados del artículo}\label{tab:paper}
\small
\begin{tabular}{lcccc}
\toprule
 & \multicolumn{2}{c}{Artículo} & \multicolumn{2}{c}{Nuestra replicación} \\
\cmidrule(lr){2-3}\cmidrule(lr){4-5}
Outcome & Coef. & EE & Coef. & EE \\
\midrule
\multicolumn{5}{l}{\textit{Panel A: habilidades cognitivas}}\\
Bayley: Cognitive & 0.250 & (0.063) & 0.250 & (0.063) \\
Bayley: Receptive language & 0.175 & (0.063) & 0.174 & (0.063) \\
Bayley: Expressive language & 0.032 & (0.062) & 0.029 & (0.062) \\
Bayley: Fine motor & 0.072 & (0.060) & 0.073 & (0.059) \\
MacArthur: Words & 0.092 & (0.064) & 0.086 & (0.064) \\
MacArthur: Complex phrases & 0.058 & (0.055) & 0.057 & (0.056) \\
\multicolumn{5}{l}{\textit{Panel B: habilidades socioemocionales}}\\
ICQ: Difficult (-) & 0.074 & (0.045) & 0.073 & (0.045) \\
ICQ: Unsociable (-) & 0.041 & (0.054) & 0.037 & (0.055) \\
ICQ: Unstoppable (-) & 0.032 & (0.054) & 0.031 & (0.054) \\
ECBQ: Inhibitory control & $-0.003$ & (0.058) & 0.000 & (0.058) \\
ECBQ: Attentional focusing & 0.069 & (0.049) & 0.072 & (0.048) \\
\multicolumn{5}{l}{\textit{Panel C: inversiones materiales}}\\
FCI: Types of play materials & 0.215 & (0.064) & 0.217 & (0.064) \\
FCI: Coloring and drawing books & $-0.133$ & (0.056) & $-0.124$ & (0.056) \\
FCI: Toys to learn movement & $-0.048$ & (0.065) & $-0.049$ & (0.065) \\
FCI: Toys to learn shapes & 0.416 & (0.088) & 0.426 & (0.088) \\
FCI: Shop-bought toys & 0.024 & (0.061) & 0.020 & (0.061) \\
\multicolumn{5}{l}{\textit{Panel D: inversiones de tiempo}}\\
FCI: Types of play activities & 0.277 & (0.050) & 0.280 & (0.051) \\
FCI: Told a story to child & 0.138 & (0.060) & 0.153 & (0.064) \\
FCI: Read to child & 0.362 & (0.062) & 0.403 & (0.068) \\
FCI: Played with toys & 0.175 & (0.060) & 0.183 & (0.060) \\
FCI: Named things to child & 0.137 & (0.048) & 0.144 & (0.048) \\
\bottomrule
\end{tabular}
\notas{Los valores del artículo corresponden a la Tabla 2 de \citet{attanasio}. Errores estándar clusterizados por municipio entre paréntesis.}
\end{table}

\section*{2. ¿Por qué corregir la inferencia con múltiples hipótesis?}

Cada test individual a un nivel $\alpha$ tiene una probabilidad $\alpha$ de rechazar una hipótesis nula verdadera (error de tipo I). Cuando se testean muchas hipótesis a la vez, la probabilidad de cometer \emph{al menos} un error de tipo I crece rápidamente con el número de tests. Si los $m$ tests fueran independientes y todas las nulas fueran ciertas, esa probabilidad sería $1-(1-\alpha)^m$: con $\alpha=0.05$, es 26\% para los 6 outcomes del panel A y 66\% para las 21 hipótesis de la tabla. Es decir, si no corregimos, es muy probable encontrar algún efecto ``significativo'' solo por azar, y el problema se agrava si después se destacan solamente los resultados significativos.

Las correcciones buscan controlar alguna medida del error conjunto. Las dos más usadas son la \textit{family-wise error rate} (FWER), que es la probabilidad de rechazar al menos una nula verdadera dentro de la familia de hipótesis, y la \textit{false discovery rate} (FDR), que es la proporción esperada de falsos rechazos entre todas las hipótesis rechazadas.

\section*{3. Corrección de los p-valores}

Definimos cada panel de la tabla como una familia de hipótesis (A: 6 hipótesis; B, C y D: 5 cada uno), porque cada panel agrupa outcomes que miden un mismo concepto. Calculamos:

\begin{enumerate}[label=\alph*.]
  \item \textbf{Bonferroni}: $p^{B}_j=\min\{1,\ m\,p_j\}$. Equivale a testear cada hipótesis a nivel $\alpha/m$ y controla la FWER.
  \item \textbf{Holm}: se ordenan los p-valores de menor a mayor, $p_{(1)}\le\dots\le p_{(m)}$, y se compara $p_{(k)}$ con $\alpha/(m-k+1)$, deteniéndose en el primer no rechazo. El p-valor ajustado es $p^{H}_{(k)}=\max_{j\le k}\min\{1,(m-j+1)\,p_{(j)}\}$. También controla la FWER.
  \item \textbf{Benjamini, Krieger y Yekutieli (BKY, 2006)}: es un procedimiento en dos etapas que controla la FDR. En la primera etapa se aplica Benjamini--Hochberg a nivel $q'=q/(1+q)$ para estimar el número de hipótesis nulas verdaderas, $\hat m_0=m-r_1$, donde $r_1$ es la cantidad de rechazos. En la segunda etapa se vuelve a aplicar Benjamini--Hochberg a nivel $q'\,m/\hat m_0$. Reportamos los \textit{sharpened q-values} de \citet{anderson}, es decir, el menor $q$ al que se rechaza cada hipótesis.
\end{enumerate}

La Tabla~\ref{tab:corr} repite la tabla del punto 1 y agrega el p-valor original y los tres p-valores corregidos.

\begin{table}[H]\centering
\caption{Efectos del tratamiento con p-valores corregidos por múltiples hipótesis}\label{tab:corr}
\small
\textbf{Panel A: habilidades cognitivas del niño}\\[0.2em]
\tabla{"../output/Tabla 2 - Panel A.tex"}\\[0.8em]
\textbf{Panel B: habilidades socioemocionales del niño}\\[0.2em]
\tabla{"../output/Tabla 2 - Panel B.tex"}\\[0.8em]
\textbf{Panel C: inversiones materiales}\\[0.2em]
\tabla{"../output/Tabla 2 - Panel C.tex"}\\[0.8em]
\textbf{Panel D: inversiones de tiempo}\\[0.2em]
\tabla{"../output/Tabla 2 - Panel D.tex"}
\notas{Ver notas de la Tabla~\ref{tab:rep}. Las estrellas corresponden al p-valor sin corregir. Las correcciones se aplican dentro de cada panel (familias de 6, 5, 5 y 5 hipótesis). BKY: \textit{sharpened q-values} de \citet{anderson}; el valor mínimo posible es 0.001.}
\end{table}

\section*{4. Similitudes y diferencias entre las estrategias}

Bonferroni y Holm controlan la misma tasa de error (FWER) y ambos son válidos bajo cualquier estructura de dependencia entre los tests, así que no requieren suponer independencia. Holm es una versión secuencial de Bonferroni: el p-valor más chico se ajusta igual que con Bonferroni (multiplicándolo por $m$), pero los siguientes se multiplican por $m-1$, $m-2$, etc. Por eso los p-valores de Holm nunca son mayores que los de Bonferroni y el procedimiento es más potente sin perder el control de la FWER. BKY, en cambio, controla la FDR, un criterio menos exigente: tolera que una proporción controlada de los rechazos sean falsos, a cambio de más poder. Su validez está demostrada bajo independencia de los tests (y en simulaciones funciona bien con dependencia positiva). Además es adaptativo: si muchas hipótesis de la familia parecen ser falsas, el umbral de la segunda etapa se relaja. Por eso, en familias con muchos efectos reales (como el panel D), los q-valores pueden ser incluso menores que los p-valores originales.

En nuestros resultados se ve justamente ese orden: BKY $\le$ Holm $\le$ Bonferroni. Los efectos más fuertes sobreviven a las tres correcciones: Bayley cognitivo y lenguaje receptivo (panel A), tipos de materiales de juego y juguetes para aprender formas (panel C), y tipos de actividades de juego, lectura, juego con juguetes y nombrar cosas (panel D), todos con p-valores corregidos menores a 0.05. Las diferencias aparecen en los resultados más débiles:
\begin{itemize}
  \item \textit{Libros para colorear} (panel C, $p=0.028$): deja de ser significativo al 5\% con Bonferroni (0.140) y con Holm (0.084, significativo solo al 10\%), pero sigue siéndolo con BKY (0.029).
  \item \textit{Contar un cuento} (panel D, $p=0.019$): con Bonferroni pasa a ser significativo solo al 10\% (0.095), mientras que con Holm (0.019) y BKY (0.005) sigue siendo significativo al 5\%.
  \item En el panel B ningún outcome es significativo antes ni después de corregir.
\end{itemize}

\section*{5. Comparación con los resultados del artículo}

El artículo reporta, para cada medida individual, el coeficiente y el error estándar sin ajustar por múltiples hipótesis. Su forma de lidiar con la multiplicidad es resumir cada panel en un factor latente (que nosotros no podemos replicar porque no tenemos esas variables). Nuestros resultados corregidos confirman sus conclusiones principales: el programa mejoró el desarrollo cognitivo (sobre todo la parte cognitiva y el lenguaje receptivo de la escala Bayley) y aumentó las inversiones de los padres, especialmente las de tiempo y algunos materiales de juego específicos. Los efectos individuales que dejan de ser robustos con las correcciones más exigentes son la caída en los libros para colorear y el aumento en contar cuentos. Los ítems socioemocionales no son significativos de manera individual ni en el artículo ni en nuestra replicación; el efecto significativo sobre el factor socioemocional que reporta el artículo surge de agregar los ítems, algo que no podemos evaluar con nuestros datos.

Las conclusiones dependen en parte de cómo se define la familia de hipótesis. Si tratáramos las 21 hipótesis como una sola familia (una opción más conservadora), con Bonferroni y Holm seguirían siendo significativos al 5\% el Bayley cognitivo, los tipos de materiales de juego, los juguetes para aprender formas, los tipos de actividades de juego, la lectura y, solo con Holm, el juego con juguetes. Con BKY los resultados prácticamente no cambian: todos los outcomes significativos al 5\% con familias por panel lo siguen siendo.

\begin{thebibliography}{9}
\bibitem[Anderson(2008)]{anderson} Anderson, M. L. (2008). Multiple inference and gender differences in the effects of early intervention: A reevaluation of the Abecedarian, Perry Preschool, and Early Training projects. \textit{Journal of the American Statistical Association}, 103(484), 1481--1495.
\bibitem[Attanasio et al.(2020)]{attanasio} Attanasio, O., Cattan, S., Fitzsimons, E., Meghir, C., y Rubio-Codina, M. (2020). Estimating the production function for human capital: results from a randomized controlled trial in Colombia. \textit{American Economic Review}, 110(1), 48--85.
\bibitem[Benjamini et al.(2006)]{bky} Benjamini, Y., Krieger, A. M., y Yekutieli, D. (2006). Adaptive linear step-up procedures that control the false discovery rate. \textit{Biometrika}, 93(3), 491--507.
\end{thebibliography}

\newpage
\appendix
\section*{Apéndice: código}
\lstinputlisting{../programs/ps2.do}

\end{document}
